\section{Future Work: Sovereign Discovery}\label{sec:future}

The proposed Sovereign DNS/Agent Naming Service would resolve an \code{ans_name}
to a versioned agent record containing endpoints, passport references,
supported protocols, jurisdictional metadata, and evidence locations.  It
would draw on DNS lessons concerning hierarchical delegation
\citep{mockapetris1987dnsconcepts,mockapetris1987dnsimplementation}, NANDA's
agent-index direction \citep{raskar2025beyonddns}, and DID/VC representations
\citep{w3c_did_core_2022,w3c_vc_data_model_2025}.  This is proposed
architecture, not implemented resolution.

A credible implementation requires more than a lookup API.  It needs namespace
governance, proof of control, authenticated record updates, revocation, key
rotation, replay protection, privacy-preserving queries, caching semantics,
conflict handling, auditability, and resilience to compromised registries.
Jurisdiction and residency assertions need issuer and evidence policies rather
than self-assertion alone.

ATrust contracts could define which identity methods, credential types, and
handshake evidence a caller requires.  DCP-AI could inform citizenship and
control-plane semantics.  \Argorix could compile these requirements and reject
a discovered record before an adapter is planned.  An executable resolver,
credential verifier, and cryptographic handshake engine would nevertheless be
new trusted components, requiring independent threat models and tests.

Empirically, future work should generate a controlled matrix: valid and invalid
records, stale keys, revoked credentials, cross-jurisdiction residency
conflicts, unavailable registries, malicious adapters, altered artifacts, and
summary/detail contradictions.  Outcome measures should include denial
correctness, false acceptance, false rejection, verification latency, evidence
size, and recovery behavior.  Prompt robustness should be studied only with an
explicitly collected, privacy-reviewed corpus.

For evidence, signed bundles and append-only transparency publication could
provide producer authentication and replacement detection, subject to key
governance and availability analysis.  This direction should not be described
as blockchain immutability unless an actual consensus/storage system is
implemented and evaluated.

\subsection{Staged roadmap}
We sequence the proposed work as a claim-bounded roadmap, in which each stage
adds a single trust property and its own evaluation rather than bundling
assurance claims.  The version labels are planning milestones, not released
guarantees.

\begin{description}
\item[v0.1] Compiled governance with offline-verifiable EvidenceBundles
(the configuration evaluated here).
\item[v0.2] Include a source digest in the EvidenceBundle so that compiled
bytecode can be tied back to its source.
\item[v0.3] Signed bundles and an append-only transparency log for producer
authentication and replacement detection.
\item[v0.4] A live provider adapter executed under explicit sandbox
constraints, with host-level telemetry.
\item[v0.5] A DID/VC resolver and credential verification, introducing the
associated trust anchors and revocation paths.
\item[v1.0] Independent, production-oriented evaluation against controlled
workloads and pre-registered outcome definitions.
\end{description}
